% Momento di una forza rispetto a un punto O: r va da O al punto di applicazione,
% theta è l'angolo fra r e F, il braccio b = r sin(theta) è la distanza fra O e
% la retta d'azione della forza.
\begin{tikzpicture}[>=Stealth, font=\small]
  % asta incernierata in O
  \filldraw[fill=richiamo!30, draw=richiamo, thick] (-0.15,-0.15) rectangle (4.6,0.15);
  \fill[richiamo!60] (0,0) circle (0.22);
  \fill (0,0) circle (1.6pt);
  \node[below left] at (-0.1,-0.15) {$O$};
  \coordinate (A) at (4.0,0);
  \fill (A) circle (1.6pt);
  \draw[<->, formula] (0,0.45) -- (4.0,0.45) node[midway, above] {$r$};
  % forza in A a 60 gradi dall'asta
  \draw[->, attenzione, very thick] (A) -- ++(60:2.2) node[right] {$\vec F$};
  \draw (A) ++(0.55,0) arc (0:60:0.55); \node at ($(A)+(30:0.8)$) {$\theta$};
  % lo stesso angolo dentro il triangolo, opposto al vertice
  \draw[medica] (A) ++(-0.55,0) arc (180:240:0.55); \node[medica] at ($(A)+(210:0.8)$) {$\theta$};
  % retta d'azione
  \draw[dashed, attenzione!70] ($(A)+(240:2.6)$) -- ($(A)+(60:2.9)$);
  % braccio: perpendicolare da O alla retta d'azione
  \coordinate (H) at ($(A)!(0,0)!($(A)+(60:1)$)$);
  \draw[medica, thick] (0,0) -- (H) node[midway, below left] {$b = r\sin\theta$};
  \draw[medica] ($(H)!0.18cm!(0,0)$) -- ++(60:0.18) -- ($(H)+(60:0.18)$);
  % verso di rotazione
  \draw[->, richiamo, thick] (40:1.5) arc (40:120:1.5);
\end{tikzpicture}
